% Einheit 2 von 4 – Terme zusammenfassen (bank/terme/e2.jsonl, zusammenbau v0.1)
%% TODO Zweigzeile Teil 1: was hier gelernt wird (Satz in Schülersprache aus dem Einheitstitel) – Einheit 2
\einheitenkopf[e2]{Einheit 2 von 4 · Terme zusammenfassen}
\zweigzeile{ab Kl. 7, je nach Buch bis Kl. 8 · P10}
\setcounter{aufgabe}{10}

%% TODO Titel: Ich-kann-Satz fehlt in der Bank (Platzhalter: Kettenname) – Nr. 11
%% TODO Anweisung: ein Satz über der ersten Teilaufgabe fehlt in der Bank – Nr. 11
\begin{aufgabe}{Gleichartige Glieder erkennen}
\begin{teile}
\teil Unterstreiche, was zusammengehört: $4y$, $7$, $9y$
\end{teile}
\end{aufgabe}

%% TODO Titel: Ich-kann-Satz fehlt in der Bank (Platzhalter: Kettenname) – Nr. 12
%% TODO Anweisung: ein Satz über der ersten Teilaufgabe fehlt in der Bank – Nr. 12
\begin{aufgabe}{Vorzahl lesen}
\begin{teile}
\teil $a$ – Vorzahl?
\end{teile}
\end{aufgabe}

%% TODO Titel: Ich-kann-Satz fehlt in der Bank (Platzhalter: Kettenname) – Nr. 13
%% TODO Anweisung: ein Satz über der ersten Teilaufgabe fehlt in der Bank – Nr. 13
\begin{aufgabe}{Vorzeichen als Teil des Gliedes}
\begin{teile}
\teil Glieder mit Vorzeichen: $5a - 2b + 7$
\end{teile}
\end{aufgabe}

%% TODO \verfahren{…}: Verfahrensüberschrift in Sachsprache fehlt in der Bank (2.3 b)
%% TODO Titel: Ich-kann-Satz fehlt in der Bank (Platzhalter: Kettenname) – Nr. 14
%% TODO Anweisung: ein Satz über der ersten Teilaufgabe fehlt in der Bank – Nr. 14
\begin{aufgabe}{Zusammenfassen}
%% TODO \rechenplatz{n}: Schrittzahl des Grundfalls fehlt in der Bank (nur bei fester Schreibform, 2.3 b)
\begin{teile}
\teil Fasse zusammen: $2x + 5x$ \leerfeld
\teil Fasse zusammen: $9x - 4x$ \leerfeld
\teil Fasse zusammen: $2x + 4x + 3x$ \leerfeld
\teil Fasse zusammen: $5x + 3 + 2x$ \leerfeld
\teil Fasse zusammen: $4a + 3b + 2a$ \leerfeld
\teil Fasse zusammen: $y + 4y$ \leerfeld
\teil Fasse zusammen: $2x - 7x$ \leerfeld
\teil Fasse zusammen: $-3x + 8x$ \leerfeld
\teil Fasse zusammen: $4x^2 + 5x + 2x^2$ \leerfeld
\teil Fasse zusammen: $1{,}5x + 3{,}5x$ \leerfeld
\teil Vereinfache den Term $8a - 2a^2 - 5a$ und berechne seinen Wert für $a = 3$. \hfill (P10 2025 GYM)
\end{teile}
\end{aufgabe}

%% TODO \verfahren{…}: Verfahrensüberschrift in Sachsprache fehlt in der Bank (2.3 b)
%% TODO Titel: Ich-kann-Satz fehlt in der Bank (Platzhalter: Kettenname) – Nr. 15
%% TODO Anweisung: ein Satz über der ersten Teilaufgabe fehlt in der Bank – Nr. 15
\begin{aufgabe}{Malnehmen}
%% TODO \rechenplatz{n}: Schrittzahl des Grundfalls fehlt in der Bank (nur bei fester Schreibform, 2.3 b)
\begin{teile}
\teil Multipliziere: $5 \cdot 2x$ \leerfeld
\teil Multipliziere: $4x \cdot 5$ \leerfeld
\teil Multipliziere: $4y \cdot 2y$ \leerfeld
\teil Multipliziere: $(-4) \cdot 5x$ \leerfeld
\teil Multipliziere: $4a \cdot 5b$ \leerfeld
\teil Multipliziere: $2 \cdot 5x \cdot 3$ \leerfeld
\teil Multipliziere: $(-6a) \cdot 2b$ \leerfeld
\end{teile}
\end{aufgabe}

%% TODO Titel: Ich-kann-Satz fehlt in der Bank (Platzhalter: Kettenname) – Nr. 16
%% TODO Anweisung: ein Satz über der ersten Teilaufgabe fehlt in der Bank – Nr. 16
\begin{aufgabe}{Zusammenfassen – Fehler finden, Begründen, Darstellungswechsel, Anwendung}
\begin{teile}
\teil Finde den Fehler und rechne richtig: \rechnung{5a + 3 + 2a &= 10a}
\teil Kann man $4x + 4y$ zusammenfassen? Begründe.
\teil Umfang des Rechtecks (Maße in cm) – Term, zusammengefasst?

\viereck[seiten={3x,2,3x,2}]{(0,0)}{(6,0)}{(6,2)}{(0,2)}
\leerfeld
\teil Ein Kiosk verkauft Brötchen zu je $p$ Euro: am Montag 45, am Dienstag 38, am Mittwoch 52 Stück. Stelle einen Term für die Einnahmen der drei Tage auf und fasse zusammen. Wie viel nimmt der Kiosk bei $p = 0{,}40$ ein?
\end{teile}
\end{aufgabe}
